% ─────────────────────────────────────────────────────────────────────────
% §1 Introduction
% ─────────────────────────────────────────────────────────────────────────
\section{Introduction}\label{sec:intro}

Large language models are increasingly deployed as the substrate for
simulated populations in computational social science, behavioral
economics, and deliberation research
\citep{park2023generative,horton2023llm,argyle2023samples}. In a growing
number of such studies the choice of underlying LLM vendor is treated as
an implementation detail --- a single commercial API is selected,
often on the basis of cost or familiarity, and the resulting simulated
population is analyzed as though its outputs were vendor-invariant. A
parallel literature on LLM alignment has separately documented that
different vendors' models exhibit systematically different refusal
behavior on politically sensitive content \citep{rottger2024safetyprompts,
durmus2024opinionqa}. What has been less systematically explored is the
\emph{joint} distribution of these two phenomena: when a simulation
study uses an LLM to generate populations engaging with a concrete
political domain, how vendor-dependent are the resulting refusal
patterns, and what does that vendor dependency look like when examined
at the appropriate resolution?

This paper audits five commercial large language models ---
\vendor{OpenAI gpt-4o-mini}, \vendor{Google gemini-2.5-flash-lite},
\vendor{xAI grok-4-fast}, \vendor{DeepSeek deepseek-chat}, and
\vendor{Moonshot kimi-k2} --- on a bank of 200 Traditional Chinese
prompts constructed to probe Taiwan political sensitivity. Each vendor
responds to each prompt under a fixed generation configuration, yielding
1{,}000 prompt--vendor observations. We hand-label each response into a
four-category refusal taxonomy (\emph{hard refusal, soft refusal,
on-task, API-blocked}), with a documented decision-tree protocol
developed and refined during labeling and released as a supplementary
artifact. Statistical analysis uses prompt-level paired bootstrap (5{,}000
resamples, BCa confidence intervals) to report uncertainty on all
primary quantities.

Four observations frame the paper's contribution. First, the intuitive
\emph{monolithic East--West alignment dichotomy} --- the framing under
which U.S.\ vendors form one behavioral cluster and Chinese-owned
vendors form another --- is empirically refuted by our data. The two
Chinese-owned vendors produce the most distant pair of refusal
distributions in the matrix (Jensen--Shannon divergence 0.200, 95\% CI
\ci{0.149}{0.256}), while \vendor{DeepSeek}'s aggregate distribution is
statistically indistinguishable from \vendor{Gemini}'s (JSD CI includes 0).
National origin does not explain the observed clustering structure.

Second, the alignment strategies we observe span two architecturally
distinct layers: a deterministic infrastructure filter operating
pre-generation, and a probabilistic RLHF-shaped model-level response
layer. Only \vendor{Kimi} operates a significant Layer 1 filter in our
dataset (rejecting 7\% of prompts with explicit filter-reason errors);
all five vendors exhibit Layer 2 refusal behavior visible in response
text. Aggregating the two layers into a single ``refusal rate''
misrepresents mechanism and flattens an important contrast among vendor
strategies.

Third, a topic-stratified analysis reveals a four-profile vendor
taxonomy that the aggregate distribution obscures. \vendor{DeepSeek} ---
aggregately similar to the Western cluster --- exhibits a
\emph{sovereignty-specific} on-task rate of 10.3\% (CI \ci{2.6}{23.3})
versus 54.0\% on non-sovereignty content, a disjoint-CI collapse unique
in our panel. \vendor{Grok} is topic-agnostic permissive (82--83\%
across topics). \vendor{Kimi}, post Layer 1 filter, is similarly
topic-agnostic permissive. The Western pair (\vendor{OpenAI},
\vendor{Gemini}) exhibits moderate uniform sovereign dampening.
These four profiles --- moderate sovereign dampening / topic-specific
RLHF collapse / topic-agnostic permissive / infra-filter-plus-permissive-model
--- organize the empirically observed vendor behavior more cleanly than
any single-axis framing we have encountered in the literature.

Fourth, \vendor{Kimi}'s pre-generation filter is best characterized
as \emph{Taiwan-statehood blocking} rather than as
\emph{sovereignty-opinion blocking}: 4 of 50 OT-expected neutral
factual prompts about specific RoC state institutions (legislature
seat counts, constitutional amendment counts, presidential term lengths,
and flag design provenance) are rejected by the filter, alongside the 9
rejected HR-expected sovereignty-opinion prompts. The filter appears to
trigger on institutional-recognition content regardless of whether the
prompt requests any opinion, a finding we believe has not been
previously documented in the vendor-audit literature.

\paragraph{Contributions.}
The paper contributes: (i) a first public audit of refusal behavior
across five major LLM vendors on a common Taiwan-political prompt bank
(§\ref{sec:method}); (ii) statistical evidence against a monolithic
East--West alignment dichotomy via pairwise JSD on refusal distributions
with disjoint 95\% CIs (§\ref{sec:results-finding1});
(iii) characterization of \vendor{Kimi}'s pre-generation filter as
Taiwan-statehood-triggered, with enumeration of all 14 blocked prompts
(§\ref{sec:results-finding2}, Appendix~\ref{app:blocked});
(iv) a two-layer architecture decomposition of vendor refusal behavior
(§\ref{sec:results-finding4}); (v) a four-profile vendor taxonomy
derived from topic-stratified on-task rates, with CI-disjoint evidence
for the \vendor{DeepSeek}--sovereignty profile
(§\ref{sec:results-finding5}); (vi) a two-tier HR$\to$SR elasticity
finding distinguishing responsive-RLHF, ceiling-bound, and
stiff-RLHF vendor regimes (§\ref{sec:results-finding7}); and (vii)
explicit recommendations for vendor-aware refusal auditing and for
multi-vendor replication in LLM agent simulation studies
(§\ref{sec:discussion}). All code, prompts, per-response logs,
hand-labels, and the auxiliary judge audit trail are released with
the paper.

\paragraph{Scope.}
We audit a five-vendor panel on a single time-slice (April 2026 model
identifiers, explicitly frozen) and a single political domain (Taiwan).
Our findings characterize those specific vendors in that specific
window; generalization across vendors outside the panel, across
political domains other than Taiwan, or across model revisions in
future time-slices is not directly supported by this audit and
requires its own replication. The paper's methodological
contribution --- the decomposition of refusal behavior into
topic~$\times$~vendor~$\times$~layer with paired bootstrap confidence
intervals --- is designed to be directly reusable for such
replications.
